\chapter{Fase 1: Preparación de la Raspberry Pi}

\section{Objetivo}

Al final de esta fase, la Raspberry Pi debe estar:

\begin{itemize}
  \item Booteando con Raspberry Pi OS Lite (64-bit).
  \item Accesible por SSH desde la Mac de trabajo.
  \item Con IP estática reservada en el router.
  \item Con Asterisk 20 instalado y respondiendo en su CLI.
  \item Con los archivos de configuración del proyecto colocados en \texttt{/etc/asterisk/}.
\end{itemize}

\section{Materiales}

\begin{itemize}
  \item Raspberry Pi 4 (2\,GB)
  \item MicroSD 32\,GB
  \item Fuente USB-C oficial
  \item Cable micro-HDMI
  \item Teclado USB (solo para el primer arranque)
  \item Televisor o monitor con HDMI
  \item Mac (u otra computadora) con lector de microSD
\end{itemize}

\section{Paso 1: Flashear la microSD}

Desde la Mac, descargar e instalar el \emph{Raspberry Pi Imager} desde el sitio oficial \texttt{raspberrypi.com/software}. Conectar la microSD a la Mac (con adaptador USB si es necesario).

Al abrir el Imager configurar:

\begin{itemize}
  \item \textbf{Device:} Raspberry Pi 4
  \item \textbf{Operating System:} Raspberry Pi OS Lite (64-bit)
  \item \textbf{Storage:} la microSD insertada
\end{itemize}

Antes de presionar \emph{Write}, hacer clic en el ícono de engranaje (\textsf{\textbf{\textcolor{hugoblue}{Edit Settings}}}) para configurar:

\begin{itemize}
  \item \textbf{Hostname:} \texttt{hugo-pbx}
  \item \textbf{Usuario:} \texttt{beto} (o el nombre que se prefiera), con contraseña robusta.
  \item \textbf{SSH:} habilitado, autenticación por contraseña o por clave pública.
  \item \textbf{WiFi:} SSID y password (solo si no se usará Ethernet).
  \item \textbf{Locale:} \texttt{America/Santiago}, teclado \texttt{latam}.
\end{itemize}

Confirmar y dejar que el Imager escriba (entre 5 y 15 minutos según la velocidad del lector).

\section{Paso 2: Primer arranque}

Insertar la microSD en la Pi, conectar HDMI, teclado USB y, finalmente, la fuente. El primer boot tarda entre 1 y 2 minutos porque expande automáticamente el sistema de archivos. Al terminar mostrará un \emph{login prompt}.

Iniciar sesión con el usuario y contraseña configurados en el paso anterior.

\section{Paso 3: Configuración de red}

Lo deseable es asignar una IP estática reservada en el router para la dirección MAC de la Pi. Para obtener la MAC del puerto Ethernet:

\begin{shellcode}
ip link show eth0 | grep ether
\end{shellcode}

La salida tendrá la forma \texttt{link/ether dc:a6:32:xx:xx:xx}. Anotar esa MAC, ingresar al panel de administración del router (habitualmente \texttt{http://192.168.1.1}) y reservar la IP \texttt{192.168.1.10} para esa MAC.

Reiniciar la Pi y verificar la IP:

\begin{shellcode}
ip addr show eth0
\end{shellcode}

\begin{tip}
Si el router asigna IPs en otro rango (por ejemplo \texttt{192.168.0.x} o \texttt{10.0.0.x}), adaptar todas las direcciones del manual en consecuencia.
\end{tip}

\section{Paso 4: Conexión por SSH}

Desde la Mac:

\begin{shellcode}
ssh beto@192.168.1.10
\end{shellcode}

Si el router resuelve mDNS, también funciona:

\begin{shellcode}
ssh beto@hugo-pbx.local
\end{shellcode}

A partir de este punto el teclado USB y el HDMI ya no son necesarios.

\section{Paso 5: Actualización del sistema}

\begin{shellcode}
sudo apt update
sudo apt upgrade -y
sudo reboot
\end{shellcode}

Reconectar SSH después del reinicio.

\section{Paso 6: Clonar el repositorio del proyecto}

Copiar el directorio \texttt{hugo-phone/} entregado en el repositorio del proyecto a la Pi. Si está en un servicio Git:

\begin{shellcode}
cd ~
git clone <URL-del-repositorio> hugo-phone
cd hugo-phone
\end{shellcode}

Si los archivos están solo en la Mac:

\begin{shellcode}
# Desde la Mac
scp -r hugo-phone/ beto@hugo-pbx.local:~/
\end{shellcode}

\section{Paso 7: Ejecutar el instalador}

\begin{shellcode}
cd ~/hugo-phone
sudo bash scripts/install-all.sh
\end{shellcode}

El script realiza, en orden:

\begin{enumerate}
  \item \texttt{apt update \&\& apt upgrade} completos.
  \item Instalación de Asterisk 20, DOSBox-Staging, \texttt{build-essential}, \texttt{libudev-dev} y demás dependencias.
  \item Copia de los archivos \texttt{sip.conf}, \texttt{extensions.conf} y \texttt{manager.conf} a \texttt{/etc/asterisk/} con permisos correctos.
  \item Instalación de \emph{rustup} en el usuario \texttt{beto} (descarga e instalación de la toolchain estable).
  \item Compilación del bridge en modo \texttt{release}, con LTO y \texttt{strip}.
  \item Creación de la regla \texttt{udev} para \texttt{/dev/uinput} y carga del módulo del kernel.
  \item Instalación del servicio \texttt{systemd} \texttt{hugo-bridge.service}.
  \item Despliegue de la configuración de DOSBox-Staging en \texttt{\~{}/.config/dosbox/}.
\end{enumerate}

La fase más lenta es la compilación del bridge, que en una Pi 4 toma entre 5 y 8 minutos. Tiempo total estimado: 20--30 minutos según velocidad de red.

\section{Paso 8: Configurar secretos}

Las contraseñas iniciales son placeholders y deben reemplazarse antes de cualquier uso real:

\begin{shellcode}
sudo nano /etc/asterisk/sip.conf
sudo nano /etc/asterisk/manager.conf
sudo nano /opt/hugo-phone/bridge/config.toml
\end{shellcode}

Reemplazar todas las ocurrencias de \texttt{cambiame-101}, \texttt{cambiame-102} y \texttt{cambiame-ami}.

\begin{advertencia}
La contraseña en \texttt{manager.conf} (sección \texttt{[hugo-bridge]}) y la del archivo \texttt{config.toml} (sección \texttt{[ami]}) deben ser \textbf{idénticas}. De lo contrario el bridge no podrá autenticarse en Asterisk.
\end{advertencia}

Aplicar cambios reiniciando Asterisk:

\begin{shellcode}
sudo systemctl restart asterisk
\end{shellcode}

\section{Paso 9: Verificación}

Ingresar a la consola interactiva de Asterisk:

\begin{shellcode}
sudo asterisk -rvvv
\end{shellcode}

Dentro del CLI ejecutar:

\begin{plaincode}
hugo-pbx*CLI> sip show peers
\end{plaincode}

Debe listar las extensiones 101 y 102 (esta última se usará en la siguiente fase) con estado \texttt{Unmonitored} (sin clientes registrados aún).

\begin{plaincode}
hugo-pbx*CLI> manager show settings
\end{plaincode}

Debe indicar \texttt{Enabled: Yes} y \texttt{Port: 5038}.

\begin{plaincode}
hugo-pbx*CLI> dialplan show hugo-game
\end{plaincode}

Debe mostrar la extensión \texttt{9000} con sus pasos de \texttt{Answer}, \texttt{Wait}, \texttt{Playback(beep)}, \texttt{Read} y \texttt{UserEvent}.

Si los tres comandos responden sin error, la \textbf{Fase 1 está completa}.

\section{Resolución de problemas}

\subsection{Asterisk no inicia}

\begin{shellcode}
sudo journalctl -u asterisk -n 80
\end{shellcode}

Casi siempre se trata de un error de sintaxis en uno de los archivos \texttt{.conf} que el log indica con número de línea exacto.

\subsection{La Pi se reinicia sin causa aparente o presenta lag aleatorio}

Probablemente la fuente de poder es insuficiente. La Pi muestra un ícono de rayo amarillo en pantalla cuando detecta \emph{undervoltage}. Cambiar a la fuente oficial 5\,V/3\,A.

\subsection{SSH no responde}

\begin{itemize}
  \item Verificar que el cable Ethernet esté conectado o que el WiFi haya tomado correctamente la red. Conectar teclado + HDMI y ejecutar \texttt{ip addr}.
  \item Revisar que el servicio SSH esté corriendo: \texttt{sudo systemctl status ssh}.
  \item Confirmar que el firewall del router no esté bloqueando el tráfico LAN.
\end{itemize}

\subsection{El instalador falla durante la compilación del bridge}

\begin{itemize}
  \item Reintentar manualmente como usuario \texttt{beto}:
  \begin{shellcode}
cd /opt/hugo-phone/bridge
~/.cargo/bin/cargo build --release
  \end{shellcode}
  \item Si falla por falta de memoria (poco probable en Pi 4 2\,GB, pero posible), agregar \emph{swap}:
  \begin{shellcode}
sudo dphys-swapfile swapoff
sudo sed -i 's/CONF_SWAPSIZE=.*/CONF_SWAPSIZE=2048/' \
  /etc/dphys-swapfile
sudo dphys-swapfile setup
sudo dphys-swapfile swapon
  \end{shellcode}
\end{itemize}
